\documentclass[11pt]{article}
\usepackage[margin=1in]{geometry}
\usepackage[T1]{fontenc}\usepackage[utf8]{inputenc}
\usepackage{lmodern,amsmath,amssymb,amsthm,mathtools,booktabs,microtype}
\usepackage[colorlinks=true,linkcolor=blue,citecolor=blue,urlcolor=blue]{hyperref}
\hypersetup{pdftitle={Agnostic learning with a fixed heavy-tailed neural-network prior},pdfauthor={ChatGPT-6 Astra; Claude Opus 5.5; Kaarel Hänni; Dmitry Vaintrob}}
\newtheorem{theorem}{Theorem}[section]\newtheorem{lemma}[theorem]{Lemma}
\newtheorem{corollary}[theorem]{Corollary}\newtheorem{proposition}[theorem]{Proposition}
\theoremstyle{remark}\newtheorem{remark}[theorem]{Remark}
\newcommand{\E}{\mathbb E}\newcommand{\R}{\mathbb R}\newcommand{\Pp}{\mathbb P}
\newcommand{\KL}{\operatorname{KL}}\newcommand{\clip}{\operatorname{clip}}
\newcommand{\whR}{\widehat R}\newcommand{\calG}{\mathcal G}
\title{Agnostic learning with a fixed heavy-tailed\\neural-network prior}
\author{ChatGPT-6 Astra \and Claude Opus 5.5 \and Kaarel H\"anni \and
Dmitry Vaintrob\thanks{Authors are listed in alphabetical order.}}
\date{17 November 2026}
\begin{document}\maketitle
\begin{abstract}
We prove an agnostic learning guarantee for one dense Bayesian neural network with independent heavy-tailed weights and a fixed Gaussian likelihood. The architecture is fixed, every coordinate has a continuous density, and there is no prior over network sizes. For a Student-family tail index $0<\alpha<2$, a row of fan-in $F$ uses weight scale $F^{-1/\alpha}$. Hidden activations are bounded with bounded first and second derivatives, and the linear readout is passed through a bounded output link, such as tanh or clipping. At every finite learner width above comparator width, the ordinary final posterior mean achieves the noise floor plus $(1+\eta)$ times the best smaller network's excess risk, with remainder of order $n^{-(2-\alpha)/2}$ for any fixed $\eta>0$. This allows arbitrary conditional sub-Gaussian label noise. An all-data expectation theorem has coefficient exactly one on comparator risk at the same sample-size rate. Independent Cauchy weights give $n^{-1/2}$; tail index $\alpha=1/2$ gives $n^{-3/4}$. Constants are independent of learner width and polynomial in comparator width, coefficient bound and input dimension at fixed depth. A complementary bounded-label theorem has leading coefficient exactly one, with a slower general-risk rate and an observable temperature calibration. The proof shows that the multiplicity of possible small-network implementations compensates the width-dependent probability of each designated large weight, while truncated background coordinates cancel in mean square.
\end{abstract}

\section{The question and the guarantee}
A useful agnostic guarantee would say that, after $n$ examples, a Bayesian neural network predicts almost as well as the best smaller network whose complexity is polynomially controlled by $n$. Such a result should continue to hold when the learner is much wider than that comparator, and when the true regression function is not itself a small network.

This paper gives such a guarantee for independent heavy-tailed coordinates. The learner has one dense architecture, fixed before seeing the observations. Every coordinate has an ordinary continuous prior density. The output is passed through a fixed bounded link; tanh and clipping are both allowed. Fixed Gaussian likelihood variances give two guarantees: an exact coefficient-one oracle in expectation, and a high-probability oracle with an arbitrarily small multiplicative factor on approximation error above the noise floor. For bounded labels, a separate calibrated theorem has coefficient one with high probability, at a slower general-risk rate.

There are two relevant features of the prior. First, a row with $N$ coordinates offers $N$ opportunities for an order-one weight, each with probability of order $1/N$. Second, after removing those large weights, the remaining centered coordinates have a controlled sum of variances. These facts produce a neighborhood of any bounded-weight comparator with prior mass bounded below independently of learner width. Standard PAC--Bayes inequalities then suffice; no separate capacity bound for the full dense network is needed.

The result is statistical. It makes no assertion that exact posterior inference, the partition function, or the temperature selection can be computed efficiently. Its dependence on teacher depth can be large. The proof uses an unrolled computation tree, whose size may be exponential in depth.

\paragraph{Related work.}
Castillo and Egels~\cite{castillo-egels} prove adaptive posterior contraction for dense ReLU networks with independent heavy-tailed weights and deterministic architectures, under prescribed architecture and shrinkage regimes. Thus dense independent heavy-tailed priors without an architecture or sparsity hyperprior are already an established approach. Our target here is a finite-sample agnostic oracle against arbitrary smaller bounded-weight networks, with constants uniform over every prescribed finite learner width above comparator width and fan-in normalization that does not depend on sample size. The model-specific step is the width-uniform function-neighborhood calculation. The distinction is not between an ordinary and a fractional Gaussian likelihood: tempering a Gaussian likelihood is equivalent to increasing its observation variance.

\section{Model and main result}
Fix input dimension $d\ge1$, hidden depth $H\ge1$, and a function $\phi\in C^2(\R)$ satisfying
\begin{equation}
 |\phi(t)|\le1,\qquad |\phi'(t)|\le A,\qquad
 |\phi''(t)|\le J,\qquad A\ge1.
 \label{eq:activation}
\end{equation}
Tanh is an example. Fix an output link $\psi:\R\to[-1,1]$, either $\clip_{[-1,1]}$ or a continuously differentiable function with $|\psi'|\le L_\psi$ and a $J_\psi$-Lipschitz derivative, where $L_\psi\ge1$. Thus tanh can be used in every layer. Put
\[
 B_\psi=\begin{cases}5,&\psi=\clip_{[-1,1]},\\
 4L_\psi+L_\psi^2+2J_\psi,&\text{for a smooth output link}.
 \end{cases}
\]
Inputs lie in $[-1,1]^d$. For the main theorem, the Borel input-label law has regression function $f_0(x)=\E[Y\mid X=x]$ in $[-1,1]$ and conditionally centered sub-Gaussian noise $\xi=Y-f_0(X)$ with proxy $\sigma^2$: $\E[e^{t\xi}\mid X]\le e^{\sigma^2t^2/2}$ for all real $t$. Write $R_0=\E\xi^2$ for the noise floor and $R(f)=\E_D(f(X)-Y)^2$. The sample $S=((X_i,Y_i))_{i=1}^n$ is IID from $D$, with $n\ge2$. The complementary calibrated theorem below assumes bounded labels $Y\in[-1,1]$; this is a special case of the main setup with $\sigma^2=1$.

For $\bar x=(1,x)$ and a fixed hidden width $N$, the dense learner is
\begin{equation}
 h^1=\phi(W^1\bar x),\qquad
 h^\ell=\phi(W^\ell h^{\ell-1})\ (2\le\ell\le H),\qquad
 z_w=v^\top h^H,\qquad f_w=\psi(z_w).
 \label{eq:network}
\end{equation}
All matrices are fully connected. Biases occur in the first layer through the constant input. Section~\ref{sec:biases} discusses other biases.

Fix $\alpha\in(0,2)$ and let
\begin{equation}
 p_\alpha(t)=c_\alpha(1+t^2)^{-(1+\alpha)/2},\qquad
 c_\alpha=\frac{\Gamma((1+\alpha)/2)}{\sqrt\pi\,\Gamma(\alpha/2)}.
 \label{eq:student}
\end{equation}
The prior $\mu_N$ makes all coordinates independent. A coordinate in a row of fan-in $F$ has law $F^{-1/\alpha}T$, where $T$ has density~\eqref{eq:student}. Thus $F=d+1$ in the first layer, and $F=N$ in internal layers and the readout. For $\alpha=1$ these are independent Cauchy weights with scales $1/(d+1)$ and $1/N$.

For $\beta>0$, the posterior and its ordinary mean are
\begin{equation}
 d\Pi_{S,\beta}(w)=Z_{S,\beta}^{-1}
 e^{-\beta\sum_{i=1}^n(f_w(X_i)-Y_i)^2}\,d\mu_N(w),
 \qquad \bar f_{S,\beta}(x)=\int f_w(x)\,d\Pi_{S,\beta}(w).
 \label{eq:posterior}
\end{equation}
The observation variance of this Gaussian likelihood is $1/(2\beta)$. At a given temperature it has no learner-width rescaling. The chosen output link is part of the mean in the likelihood, not merely a transformation applied after inference.

Let $\calG_{m,B}$ be the class of networks with the same depth, activation, first-layer bias convention and output link, hidden width $m$, and all coefficients of absolute value at most $B\ge1$. We require $m\le N$. Set
\begin{equation}
 q=\frac\alpha{2-\alpha},\quad M=(m+1)^H,\quad
 S_{m,B}=C_\alpha(1+A^2+J)(d+1)(1+AmB)^{2H},
 \label{eq:resources}
\end{equation}
where the constant is sufficiently large that $S_{m,B}\ge1$, and define the explicit upper resource
\begin{align}
 K_{m,B}=C_\alpha M\bigg[&(B_\psi S_{m,B})^q+1\\[-2mm]
 &+(d+1)\left\{\log\!\big(C_\alpha M(d+1)(B+6)^{1+\alpha}\big)
 +\frac{\log(B_\psi S_{m,B})+1/q}{2-\alpha}+1\right\}\bigg].
 \label{eq:complexity}
\end{align}
Increase constants so that $K_{m,B}\ge1$. At fixed $H,\alpha,A,J$, this is polynomial in $d,m,B$, up to logarithms.

\begin{theorem}[A fixed-likelihood oracle above the noise floor]\label{thm:fast-oracle}
Fix the learner width $N$ before drawing the sample, comparator width $m\le N$, coefficient cap $B\ge1$, confidence $0<\delta<1$, and approximation factor $0<\eta\le1$. Use the single fixed temperature
\[
 \beta=\frac{\eta}{32(1+\sigma^2)},
 \qquad \text{Gaussian observation variance }\frac1{2\beta}
                   =\frac{16(1+\sigma^2)}\eta.
\]
Put $A_* =\min_{g\in\calG_{m,B}}\{R(g)-R_0\}$. With probability at least $1-\delta$,
\begin{align}
 R(\bar f_{S,\beta})
 &\le\int R(f_w)\,d\Pi_{S,\beta}(w)\\
 &\le R_0+(1+\eta)A_*
 +C_\alpha\left(
 \frac{(1+\sigma^2)\{K_{m,B}+\log(2/\delta)\}}{\eta n}
 \right)^{\frac{2-\alpha}{2}}.
 \label{eq:fast-oracle}
\end{align}
The prior and likelihood are independent of $n,m,B$ and the unknown regression function. The constants are uniform over every finite $N\ge m$ and may depend on the fixed output link.
\end{theorem}

\begin{corollary}[A simultaneous structural oracle]\label{cor:structural-oracle}
The same fixed prior and fixed likelihood satisfy~\eqref{eq:fast-oracle} simultaneously for every integer $1\le m\le N$ and integer cap $B\ge1$, after replacing $\log(2/\delta)$ by $\log(CmB/\delta)$ and increasing the constant. Consequently one may take the infimum of these bounds over all such comparator classes.
\end{corollary}
\begin{proof}
Apply Theorem~\ref{thm:fast-oracle} with failure probabilities proportional to $\delta/(m^2B^2)$ and take a union bound. This is only a proof allocation: the learner has the same prior, likelihood and posterior for every comparator class.
\end{proof}

The factor $1+\eta$ multiplies approximation error above the noise floor, not irreducible label noise. In particular, if the conditional mean itself is a smaller bounded-weight network, the ordinary final posterior has excess risk of order $n^{-(2-\alpha)/2}$ even with nonzero label noise. Choosing a smaller fixed $\eta$ makes the factor as close to one as desired, at the displayed cost in sample size.

\begin{theorem}[A sharp expected-risk oracle using all observations]
\label{thm:expected-oracle}
Keep the same prior and use the fixed temperature
$\beta_*=1/[256(1+\sigma^2)]$, with Gaussian observation variance
$128(1+\sigma^2)$. For every fixed comparator $g\in\calG_{m,B}$,
\[
 \E_{S_n}R(\bar f_{S_n,\beta_*})
 \le R(g)+C_\alpha
 \left(\frac{(1+\sigma^2)K_{m,B}}{n+1}\right)^{(2-\alpha)/2}.
\]
The constants are uniform over all finite $N\ge m$. One may take the infimum over deterministic comparators and resource classes directly. This is the ordinary final posterior trained on all observations; the guarantee is in expectation over the sample.
\end{theorem}
\begin{proof}
For each $0<r\le1$, Lemma~\ref{lem:comparison} applied with bounded reference label $f_0(x)$ gives a data-independent law with risk at most $R(g)+r$ and KL at most $K_{m,B}r^{-q}$. The generic all-data expectation theorem proved in Section~\ref{sec:sharp-expectation-proof} gives
\[
 \E R(\bar f_{S_n,\beta_*})
 \le R(g)+r+\frac{256(1+\sigma^2)K_{m,B}r^{-q}}{n+1}.
\]
Choose $r$ of order $\{(1+\sigma^2)K_{m,B}/(n+1)\}^{1/(q+1)}$ when this is at most one. Otherwise the deterministic excess-risk bound $R(f)-R_0\le4$ proves the conclusion after enlarging the constant.
\end{proof}

The exact coefficient one in this expectation theorem does not replace the high-probability scope of Theorem~\ref{thm:fast-oracle}. Both use a fixed physical likelihood variance and allow conditionally sub-Gaussian labels.

\paragraph{A complementary coefficient-one guarantee.}
For bounded labels, a different calibration gives leading coefficient exactly one, while retaining the same faster rate when comparator risk vanishes. Its general-risk rate is slower. Define the observable normalized free energy

\[
 F_{S,\beta}=-\frac1{\beta n}\log Z_{S,\beta},
\]
and the grid
\[
 \mathcal B_n=\{2^{-j}:3\le j\le\lceil\log_2(64n)\rceil\},
 \qquad c_\beta=\frac{1-e^{-4\beta}}{4\beta}.
\]
Given $0<\delta<1$, let $J_n=|\mathcal B_n|$ and select any minimizer
\begin{equation}
 \widehat\beta\in\arg\min_{\beta\in\mathcal B_n}
 \frac{F_{S,\beta}+\log(2J_n/\delta)/(\beta n)}{c_\beta}.
 \label{eq:calibration}
\end{equation}
Ties can be broken by the grid order.

\begin{theorem}[A width-uniform agnostic oracle]\label{thm:oracle}
Assume $Y\in[-1,1]$. Fix the learner width $N$ before drawing the sample, and fix $m\le N$, $B\ge1$. Put $R_* =\min_{g\in\calG_{m,B}}R(g)$ and $T_{m,B}=K_{m,B}+\log(2J_n/\delta)$. For the prior~\eqref{eq:student} and posterior~\eqref{eq:posterior}, with probability at least $1-\delta$,
\begin{align}
 R(\bar f_{S,\widehat\beta})
 &\le\int R(f_w)\,d\Pi_{S,\widehat\beta}(w)\\
 &\le R_*+C_\alpha\left[
 \left(\frac{T_{m,B}R_*}{n}\right)^{\frac{2-\alpha}{4-\alpha}}
 +\left(\frac{T_{m,B}}n\right)^{\frac{2-\alpha}2}
 \right].
 \label{eq:main-oracle}
\end{align}
The constants contain no learner-width factor. The same calibration rule applies at every fixed width, using that learner's partition function; it does not use $m,B$, or the unknown law $D$.
\end{theorem}

The comparator minima exist: the parameter cube is compact, prediction is bounded and continuous in the coefficients, and the label law has a finite second moment. The probability statement holds for each finite architecture fixed before the data; it does not assert one common event for unrestricted data-dependent selection of infinitely many learner widths.

For Cauchy weights, suppressing logarithms and fixing $H,A,J$, the resource can be taken of order
\begin{equation}
 K_{m,B}=\widetilde O_H\!\left((d+1)(m+1)^H(1+AmB)^{2H}\right).
 \label{eq:cauchy-poly}
\end{equation}
For the fixed-likelihood theorem, $n\gtrsim (1+\sigma^2)K_{m,B}/(\eta\varepsilon^2)$ suffices for remainder $\varepsilon$. For the coefficient-one bounded-label theorem, $n\gtrsim K_{m,B}(R_*\varepsilon^{-3}+\varepsilon^{-2})$ suffices to compete within $\varepsilon$ with the best comparator. At bounded coefficients this is a width degree $3H$, up to logarithms. The theorem does not claim this dependence is optimal. The oracle statements give the following rates (logarithms and confidence terms are suppressed). The first two rate columns allow conditional sub-Gaussian labels and use fixed temperatures; the last assumes bounded labels and uses the specified calibration:
\begin{center}
\begin{tabular}{cccc}\toprule
Tail index & \shortstack{Expectation\\factor one} & \shortstack{High probability\\factor $1+\eta$} & \shortstack{High probability\\factor one}\\\midrule
$\alpha=1$ (Cauchy)&$n^{-1/2}$&$n^{-1/2}$&$n^{-1/3}$\\
$\alpha=1/2$&$n^{-3/4}$&$n^{-3/4}$&$n^{-3/7}$\\
$0<\alpha<2$&$n^{-(2-\alpha)/2}$&$n^{-(2-\alpha)/2}$&$n^{-(2-\alpha)/(4-\alpha)}$\\\bottomrule
\end{tabular}
\end{center}
The dependence of $K$ and all constants on $\alpha$ remains relevant when comparing priors. We fix $\alpha$; no uniform assertion is made as $\alpha\downarrow0$ or $\alpha\uparrow2$.

\begin{corollary}[An explicit exponential depth upper bound]\label{cor:depth-upper}
Fix $\alpha,A,J,d,m,B,\eta,\sigma^2$, a positive remainder accuracy $\varepsilon$, and confidence $1-\delta$. There is a constant $C$ depending only on these fixed quantities such that, at every hidden depth $H\ge1$ and learner width $N\ge m$, at most $e^{C(H+1)}$ samples suffice for the remainder in~\eqref{eq:fast-oracle} to be at most $\varepsilon$.
\end{corollary}
\begin{proof}
The constant in $S_{m,B}$ is independent of $H$, as the moment induction below shows. The explicit resource~\eqref{eq:complexity} is therefore bounded by $e^{C(H+1)}$ when the other quantities are fixed. Theorem~\ref{thm:fast-oracle} then gives the assertion.
\end{proof}


\section{A comparator neighborhood with width-independent cost}\label{sec:comparison}
The main construction is a comparison measure, absolutely continuous with respect to the original product prior. For labels in $[-1,1]$, it will have a controlled expected empirical risk for \emph{every} sample, not only for a high-probability set of samples. The noisy-regression proof applies this bounded-label statement to the reference label $f_0(x)$ and then concentrates the comparison law's relative loss.

\subsection{Small coordinates and isolated large coordinates}
\begin{lemma}[Truncated Student coordinates]\label{lem:truncation}
Let $W=F^{-1/\alpha}T$, with $F\ge1$ and $T$ having density~\eqref{eq:student}. For $0<\epsilon\le1$,
\begin{align}
 -F\log\Pp(|W|\le\epsilon)&\le C_\alpha\epsilon^{-\alpha},\\
 \E[W\mid |W|\le\epsilon]&=0,\\
 F\E[W^2\mid |W|\le\epsilon]&\le C_\alpha\epsilon^{2-\alpha}.
 \label{eq:truncation}
\end{align}
For $|w|\le B+5$, the unconditioned density satisfies
\begin{equation}
 p_{F,\alpha}(w)\ge\frac{c'_\alpha}{F(B+6)^{1+\alpha}}.
 \label{eq:density-lower}
\end{equation}
\end{lemma}
\begin{proof}
Put $u=F^{1/\alpha}\epsilon$. For $0<u\le1$, the small-ball probability is at least $c'_\alpha u$. Since $F=u^\alpha\epsilon^{-\alpha}$ and $u^\alpha\log(1/(c'_\alpha u))$ is bounded, the first inequality follows in this range. For $u\ge1$, the tail probability is at most $C_\alpha u^{-\alpha}$ and the small-ball probability is bounded below by a positive constant. Consequently $-\log\Pp(|T|\le u)\le C_\alpha u^{-\alpha}$, proving the other range.

Symmetry gives the conditional mean. For $u\le1$, the conditional second moment is at most $\epsilon^2$, and $F\epsilon^2\le\epsilon^{2-\alpha}$. For $u\ge1$, direct integration gives $\E[T^2\mathbf1_{|T|\le u}]\le C_\alpha u^{2-\alpha}$. Divide by the small-ball probability and rescale to obtain~\eqref{eq:truncation}. Finally,
\[
 p_{F,\alpha}(w)=\frac{c_\alpha}{F}
 (F^{-2/\alpha}+w^2)^{-(1+\alpha)/2},
\]
which implies~\eqref{eq:density-lower}.
\end{proof}

We need designated teacher edges to be distinguishable from the truncated background, even when the teacher coefficient vanishes. Deleting small teacher edges would introduce a first-order bias. The following unbiased construction avoids that loss.

\begin{lemma}[Unbiased designated edges]\label{lem:active}
For every $a\in[-B,B]$ and $0<\epsilon\le1$ there is an absolutely continuous law $Q_a$ such that
\begin{equation}
 \E_{Q_a}W=a,\quad \operatorname{Var}_{Q_a}(W)\le17\epsilon^2,
 \quad Q_a\{|W|>\epsilon\}=1,
 \label{eq:active-properties}
\end{equation}
whose density is bounded by $1/(2\epsilon)$ and whose support lies in $[-B-5,B+5]$. For the fan-in-$F$ prior,
\begin{equation}
 \KL(Q_a\|p_{F,\alpha})\le
 \log\frac{C_\alpha F(B+6)^{1+\alpha}}\epsilon.
 \label{eq:active-kl}
\end{equation}
\end{lemma}
\begin{proof}
If $|a|>4\epsilon$, use the uniform law on $[a-\epsilon,a+\epsilon]$. Otherwise mix the uniform laws on $[3\epsilon,5\epsilon]$ and $[-5\epsilon,-3\epsilon]$ with probabilities $(1+a/(4\epsilon))/2$ and $(1-a/(4\epsilon))/2$. The mixture has mean $a$ and variance at most $(16+1/3)\epsilon^2<17\epsilon^2$. The density and support claims are immediate, and~\eqref{eq:active-kl} follows from~\eqref{eq:density-lower}. An ordinary symmetric interval around $a$ has the same KL upper bound, without the separated-support property.
\end{proof}

\subsection{Unrolling and counting implementations}
A width-$m$, depth-$H$ network can be unrolled from its scalar readout into a formula: duplicate a hidden neuron whenever its value is used by different parents. The original function is unchanged. The formula has at most
\[
 1+m+\cdots+m^H\le M=(m+1)^H
\]
nodes including the scalar root. Each non-root node has exactly one parent. This last property is what allows implementation multiplicity to compensate the probability of a large edge.

Suppose first that $N\ge M$. In each hidden layer, assign every formula-node occurrence a disjoint group of physical indices of size at least $N/(2M)$. Choose one index independently and uniformly from each group. Conditional on this assignment, construct a product comparison law as follows:
\begin{itemize}
\item A selected first-layer row has independent uniform coordinates of radius $\epsilon$ around the corresponding teacher row.
\item In each selected later row, including the scalar readout, designated child edges have the laws $Q_a$ from Lemma~\ref{lem:active}; every other incoming coordinate has its prior law conditioned on $|W|\le\epsilon$.
\item All unselected rows retain their unconditioned product prior.
\end{itemize}
Every coordinate law is absolutely continuous with respect to its prior. The selection of a child's index depends only on the independent assignment variable, not on the incoming weights that compute that child's value.

\begin{lemma}[Disjoint implementation supports]\label{lem:disjoint}
Distinct complete index assignments give mutually singular comparison components. If $\nu$ is their uniform mixture, then
\[
 \KL(\nu\|\mu_N)=\text{average component KL}
                 -\text{entropy of the index assignment}.
\]
\end{lemma}
\begin{proof}
The scalar readout has, in each designated child group, exactly one coordinate with magnitude greater than $\epsilon$. It therefore determines the chosen indices in the last hidden layer. The incoming rows at those indices determine the choices one layer earlier by the same rule, and so on. Thus a weight vector in the support decodes a unique complete assignment. On a component support the mixture density is its component density times the assignment probability. Integrating its log density ratio gives the identity.
\end{proof}

For each designated edge, the active-coordinate cost~\eqref{eq:active-kl} contains $\log N$. Its child-index entropy is at least $\log(N/(2M))$. These terms cancel, leaving at most $\log(2M)$. There is one such cancellation for every non-root formula node. A selected row's background costs at most $C_\alpha\epsilon^{-\alpha}$ by Lemma~\ref{lem:truncation}. First-layer designated coordinates incur their ordinary interval costs, with fan-in $d+1$.

If instead $m\le N<M$, place the original width-$m$ teacher at fixed physical indices. Use the same designated and background coordinate laws, without unrolling or averaging over assignments. Now $\log N\le\log M$ already. There are at most $Hm+1\le M$ selected rows, at most $(H-1)m^2+m\le M$ designated internal and readout edges, and at most $M(d+1)$ designated first-layer coordinates.

Both constructions prove the following bound, uniformly over all $N\ge m$:
\begin{equation}
 \KL(\nu_{g,\epsilon,N}\|\mu_N)
 \le C_\alpha M\left[
 \epsilon^{-\alpha}+(d+1)
 \log\frac{C_\alpha M(d+1)(B+6)^{1+\alpha}}\epsilon
 \right].
 \label{eq:embedding-kl}
\end{equation}
The formula, its index groups and the comparison mixture are proof devices. The learner's prior remains the single dense independent-coordinate prior~\eqref{eq:student}.

\section{Second-order risk control}
The designated coefficient perturbations are unbiased. Together with bounded hidden features, this gives a small variance and a second-order mean bias even though the unconditioned prior has infinite coordinate variance.

\begin{lemma}[Pointwise perturbation moments]\label{lem:perturbation}
Let $g\in\calG_{m,B}$, and write $z_g$ for its raw linear readout. The comparison measure in~\eqref{eq:embedding-kl} satisfies
\begin{equation}
 \sup_x\E_\nu|z_w(x)-z_g(x)|^2
 \le S_{m,B}\epsilon^{2-\alpha},\qquad
 \sup_x\left|\E_\nu[z_w(x)-z_g(x)]\right|
 \le S_{m,B}\epsilon^{2-\alpha}.
 \label{eq:raw-moments}
\end{equation}
Here $S_{m,B}$ may be taken as in~\eqref{eq:resources}, independently of learner width.
\end{lemma}
\begin{proof}
First condition on a particular index assignment in the large-width construction; the small-width construction is already a product law. All selected incoming rows are independent of earlier layers. All hidden coordinates, selected or unselected, have magnitude at most one.

Put $\theta=\epsilon^{2-\alpha}$; then $\epsilon^2\le\theta$. Let $v_\ell$ be the largest second moment of a selected hidden neuron's error relative to its corresponding teacher value, and let $b_\ell$ be the largest absolute mean error. We fix an input throughout; all bounds are uniform over inputs and assignments.

In the first layer the preactivation error is centered and has second moment at most $(d+1)\epsilon^2/3$. At a later selected node, write the preactivation error as
\[
 \Delta=\sum_{j=1}^m a_j(h_j-h_j^*)+Z,
\]
where $a_j$ are the teacher coefficients, $h_j^*$ the teacher features, and $Z$ contains the designated coefficient perturbations and the background sum. Conditional on all preceding layers,
\begin{equation}
 \E[Z\mid\text{preceding layers}]=0,\qquad
 \E[Z^2\mid\text{preceding layers}]
 \le C_\alpha(m+1)\theta.
 \label{eq:centered-background}
\end{equation}
Indeed, the incoming coordinates are independent and centered after subtraction of their designated means. Their variances are bounded by Lemmas~\ref{lem:truncation} and~\ref{lem:active}, and all multiplying hidden features are bounded by one. This argument does not require the unselected hidden neurons to be close to zero.

The conditional centering removes the cross term, so
\[
 \E\Delta^2\le m^2B^2v_{\ell-1}+C_\alpha(m+1)\theta,
 \qquad |\E\Delta|\le mB b_{\ell-1}.
\]
Lipschitzness and Taylor's theorem therefore give
\begin{align}
 v_\ell&\le A^2\{m^2B^2v_{\ell-1}+C_\alpha(m+1)\theta\},\\
 b_\ell&\le AmB b_{\ell-1}
       +\frac J2\{m^2B^2v_{\ell-1}+C_\alpha(m+1)\theta\}.
 \label{eq:moment-recurrence}
\end{align}
The initial bounds are $v_1\le A^2(d+1)\theta/3$ and $b_1\le J(d+1)\theta/6$.

For completeness, set $a=AmB\ge1$ and $G=1+a$. Since $m+1\le2a$ and $G^2-a^2=1+2a$, induction in the first recurrence gives
\[
 v_\ell\le C_\alpha A^2(d+1)G^{2\ell-2}\theta.
\]
The second recurrence is absorbed by $G^2-a\ge a^2+a+1$, giving
\[
 b_\ell\le C_\alpha J(d+1)G^{2\ell-2}\theta.
\]
The same constants work at every depth: the gaps $G^2-a^2$ and $G^2-a$ absorb each additive term in the displayed recurrences, so no layer-dependent prefactor is introduced. The readout has the same conditional variance estimate, but no activation. Its second moment is at most $m^2B^2v_H+C_\alpha(m+1)\theta$, and its mean error is at most $mB b_H$. Both are bounded by a constant times $(1+A^2+J)(d+1)G^{2H}\theta$. This proves~\eqref{eq:raw-moments} for every assignment and hence also for their mixture.
\end{proof}

A Lipschitz estimate on the loss would only turn a mean-square error $\theta$ into a risk error of order $\sqrt\theta$. The next property preserves the second-order improvement, including at clipping thresholds.

\begin{lemma}[Squared loss after a bounded output link]\label{lem:semiconcavity}
For $|y|\le1$, let $\ell_y(z)=(\psi(z)-y)^2$. At every $z\in\R$ there is $d_y(z)$ such that, for every $u\in\R$,
\begin{equation}
 \ell_y(z+u)\le\ell_y(z)+d_y(z)u+C_\psi u^2,
 \qquad |d_y(z)|\le D_\psi,\qquad D_\psi+C_\psi\le B_\psi.
 \label{eq:semiconcavity}
\end{equation}
For clipping one can use $D_\psi=4$, $C_\psi=1$. For the smooth link class, one can use $D_\psi=4L_\psi$ and $C_\psi=L_\psi^2+2J_\psi$.
\end{lemma}
\begin{proof}
For clipping, $\ell_y(z)-z^2$ is concave: it is affine inside $[-1,1]$, has second derivative $-2$ outside, and its derivative jumps downward at both endpoints. Its supergradient inequality proves the claim, with $d_y(z)=2(z-y)$ for $|z|<1$ and zero for $|z|\ge1$.
For a smooth link, $|\ell_y'|\le4L_\psi$. Its derivative is Lipschitz with constant at most $2L_\psi^2+4J_\psi$, because $|\psi-y|\le2$. Taylor's inequality proves the displayed estimate.
\end{proof}

\begin{lemma}[A deterministic comparison for every sample]\label{lem:comparison}
For every $g\in\calG_{m,B}$, $0<r\le1$ and $N\ge m$, there is a measure $\nu_{g,r,N}\ll\mu_N$, independent of the training sample, satisfying
\begin{equation}
 \E_\nu(f_w(x)-y)^2\le(g(x)-y)^2+r
 \quad\text{for every }x\in[-1,1]^d,\ |y|\le1,
 \qquad \KL(\nu\|\mu_N)\le K_{m,B}r^{-q}.
 \label{eq:comparison}
\end{equation}
In particular the empirical-risk inequality holds for every sample with labels in $[-1,1]$, and the population-risk inequality holds for every law with labels in that interval. Section~\ref{sec:noisy-regression} uses it with the bounded reference label $f_0(x)$ to treat unbounded observed labels.
\end{lemma}
\begin{proof}
By Lemmas~\ref{lem:perturbation} and~\ref{lem:semiconcavity},
\[
 \E_\nu\ell_y(z_w(x))\le\ell_y(z_g(x))+B_\psi S_{m,B}\epsilon^{2-\alpha}.
\]
Choose $\epsilon=(r/(B_\psi S_{m,B}))^{1/(2-\alpha)}\le1$. Then $\epsilon^{-\alpha}=(B_\psi S_{m,B})^q r^{-q}$. Substitute this in~\eqref{eq:embedding-kl}. The additional logarithm of $1/r$ is bounded using $r^q\log(1/r)\le1/(eq)$, and the remaining terms are bounded by~\eqref{eq:complexity}.
\end{proof}

\begin{corollary}[Function-neighborhood mass]\label{cor:mass}
For any input law $D_X$, any $g\in\calG_{m,B}$, and $0<t\le1$, the same prior satisfies
\begin{equation}
 \mu_N\{\|f_w-g\|_{L^2(D_X)}\le t\}
 \ge \exp\{-C_\alpha K_{m,B}t^{-2q}\},
 \qquad N\ge m.
 \label{eq:function-mass}
\end{equation}
Thus the comparison argument also supplies an actual width-independent lower bound on a function neighborhood's prior probability.
\end{corollary}
\begin{proof}
Use Lemma~\ref{lem:comparison} with the noiseless law $Y=g(X)$ and $r=t^2/2$. Markov's inequality gives comparison probability at least $1/2$ to the displayed event. The binary relative-entropy inequality gives $\KL(\nu\|\mu_N)\ge\tfrac12\log(1/\mu_N(E))-\log2$, proving the claim after increasing the constant.
\end{proof}

\section{A fixed-temperature oracle for noisy regression}
\label{sec:noisy-regression}
We now prove Theorem~\ref{thm:fast-oracle}. The central quantity is the excess risk above the regression function, which removes the irreducible noise from the complexity penalty.

Write $f_0(x)=\E[Y\mid X=x]$, $\xi=Y-f_0(X)$, and
\[
 R_0=\E\xi^2,\qquad
 A(f)=\E(f(X)-f_0(X))^2=R(f)-R_0.
\]
The assumption is $|f_0|\le1$ and
\begin{equation}
 \E[e^{t\xi}\mid X]\le e^{\sigma^2t^2/2}
 \quad\text{for every }t\in\R.
 \label{eq:conditional-subgaussian}
\end{equation}
No Gaussian assumption is needed. Bounded labels in $[-1,1]$ satisfy this with $\sigma^2=1$ by conditional Hoeffding. Gaussian label noise of variance at most $\sigma^2$ is another example.

Let $s=1+\sigma^2$ and $\beta=\eta/(32s)$, where $0<\eta\le1$. Every network output $f_w$ remains bounded, although the labels need not be. The posterior is still~\eqref{eq:posterior}; its Gaussian observation variance is the constant $16s/\eta$, independent of sample size and learner width.

\begin{lemma}[A relative-loss PAC--Bayes inequality]\label{lem:relative-pac}
With probability at least $1-\delta/2$, simultaneously for every probability measure $\rho$ on weights,
\begin{equation}
 c_\beta^{\mathrm{rel}}\E_\rho A(f_w)
 \le\E_\rho[\whR_S(f_w)-\whR_S(f_0)]
       +\frac{\KL(\rho\|\mu_N)+\log(2/\delta)}{\beta n},
 \label{eq:relative-pac}
\end{equation}
where
\[
 c_\beta^{\mathrm{rel}}=\frac{1-e^{-4\beta(1-2\beta\sigma^2)}}{4\beta}
       \ge1-4s\beta=1-\eta/8.
\]
\end{lemma}
\begin{proof}
Fix $f=f_w$ and put $e=f-f_0$. The difference of squared losses is $e(X)^2-2\xi e(X)$. Conditional sub-Gaussianity gives
\[
 \E[e^{-\beta\{(f(X)-Y)^2-(f_0(X)-Y)^2\}}\mid X]
 \le e^{-\beta(1-2\beta\sigma^2)e(X)^2}.
\]
Since $0\le e^2\le4$, convexity on that interval yields
\[
 \E e^{-\beta\{(f(X)-Y)^2-(f_0(X)-Y)^2\}}
 \le e^{-\beta c_\beta^{\mathrm{rel}} A(f)}.
\]
The lower bound on $c_\beta^{\mathrm{rel}}$ follows from $1-e^{-u}\ge u-u^2/2$ and the chosen range of $\beta$. Independence, integration against the prior, Markov's inequality and the entropy variational inequality prove~\eqref{eq:relative-pac}. The individual squared losses have finite expectation under~\eqref{eq:conditional-subgaussian}; their difference has the exponential moment just established.
\end{proof}

The comparison measure may itself be averaged before applying concentration. This is useful because the perturbation lemma is naturally a statement about its expected risk.

\begin{lemma}[Concentration for a fixed comparison law]\label{lem:comparison-concentration}
Let $\nu$ be any fixed probability measure on bounded predictions, independent of the data, and let $a=\E_\nu A(f_w)$. With probability at least $1-e^{-\ell}$,
\begin{equation}
 \E_\nu[\whR_S(f_w)-\whR_S(f_0)]
 \le a+C\sqrt{sa\ell/n}+Cs\ell/n.
 \label{eq:comparison-bernstein}
\end{equation}
\end{lemma}
\begin{proof}
Set
\[
 v(x)=\E_\nu(f_w(x)-f_0(x))^2,\qquad
 b(x)=\E_\nu(f_w(x)-f_0(x)).
\]
Then $0\le v\le4$, $b^2\le v$, and one comparison-averaged excess loss is $L=v(X)-2\xi b(X)$, of mean $a$. For $0\le t\le1/(16s)$, conditional sub-Gaussianity gives
\[
 \E e^{t(L-a)}
 \le e^{-ta}\E e^{(t+2\sigma^2t^2)v(X)}
 \le\exp\left\{a\left[\frac{e^{4(t+2\sigma^2t^2)}-1}{4}-t\right]\right\}
 \le e^{6sa t^2}.
\]
For the last inequality, put $u=t+2\sigma^2t^2\le9t/8$, use $4u\le1$ and $e^{4u}-1\le4u+16u^2$. Chernoff's bound for the independent observations then gives~\eqref{eq:comparison-bernstein}, with an absolute constant. Equivalently, optimize over $t\le1/(16s)$ to obtain a square-root term and a linear term.
\end{proof}

\begin{proof}[Proof of Theorem~\ref{thm:fast-oracle}]
Fix a deterministic risk-minimizing $g\in\calG_{m,B}$ and write $A_*=A(g)$. For a precision $0<r\le1$, Lemma~\ref{lem:comparison}, applied pointwise with label $y=f_0(x)$, supplies a data-independent comparison law with
\[
 a=\E_\nu A(f_w)\le A_*+r,
 \qquad \KL(\nu\|\mu_N)\le K_{m,B}r^{-q}.
\]
This step does not require the observed labels to be bounded.

At the Gibbs posterior, the right side of~\eqref{eq:relative-pac} is
$F_{S,\beta}-\whR_S(f_0)+\log(2/\delta)/(\beta n)$.
The free-energy variational formula bounds it by
\[
 \E_\nu[\whR_S(f_w)-\whR_S(f_0)]
       +\frac{K_{m,B}r^{-q}+\log(2/\delta)}{\beta n}.
\]
Use Lemma~\ref{lem:comparison-concentration} with $\ell=\log(2/\delta)$ and Young's inequality to obtain, on the intersection of the two events,
\[
 c_\beta^{\mathrm{rel}}\E_{\Pi_{S,\beta}}A(f_w)
 \le(1+\eta/8)(A_*+r)
       +\frac{Cs\ell}{\eta n}
       +\frac{32sK_{m,B}r^{-q}}{\eta n}.
\]
Since $c_\beta^{\mathrm{rel}}\ge1-\eta/8$, the coefficient of $A_*$ after division is at most $1+\eta$. Let
\[
 u=\left(\frac{s\{K_{m,B}+\log(2/\delta)\}}{\eta n}\right)^{1/(q+1)}.
\]
If $u\le1$, take $r=u$; every remaining term is at most $C_q u$. If $u>1$, the deterministic bound $A(f_w)\le4$ gives the result after increasing $C_q$. Finally $R(f_w)=R_0+A(f_w)$, and Jensen's inequality transfers the posterior-average bound to the ordinary final posterior mean.
\end{proof}

The unknown regression function appears only inside the proof. The learner uses neither $f_0$, $R_0$, nor the comparator. The prescribed likelihood variance depends only on the chosen approximation factor $1+\eta$ and a known conditional sub-Gaussian proxy $\sigma^2$; an upper bound on that proxy is sufficient.

% Independent proof module, 18 November 2026.
% This file is a review artifact, not an edit to an owner manuscript.
% The leave-one-out strategy is motivated by Hogsgaard--Rebeschini--Wegel,
% arXiv:2607.02247 (2026). Their stated theorem is for finite dictionaries
% and bounded losses. The arbitrary-prior and conditional sub-Gaussian
% extension below is derived here; it should receive independent checking.

\section{A sharp expectation oracle for the ordinary final posterior}
\label{sec:sharp-expectation-proof}

Let $\pi$ be any probability prior on jointly measurable functions
$f_w:\mathcal X\to[-1,1]$. Suppose the observations are IID,
$f_0(x)=\mathbb E[Y\mid X=x]\in[-1,1]$, and
$\xi=Y-f_0(X)$ is conditionally sub-Gaussian with proxy $\sigma^2$.
Set $s=1+\sigma^2$ and fix $T\ge256s$.
For a sample $S_n$ of size $n$, let
\[
 dQ_{S_n}(w)=Z_{S_n}^{-1}
   \exp\left\{-\frac1T\sum_{i=1}^n(f_w(X_i)-Y_i)^2\right\}d\pi(w),
 \qquad \bar f_{S_n}=\int f_w\,dQ_{S_n}(w).
\]
This is the ordinary posterior with Gaussian observation variance $T/2$,
using all $n$ observations. No averaging over sample sizes is performed.

\begin{theorem}[Sharp all-data expectation bound]\label{thm:generic-expectation}
For every deterministic comparison law $\rho\ll\pi$,
\[
 \mathbb E_{S_n} R(\bar f_{S_n})
 \le \int R(f_w)\,d\rho(w)
        +\frac{T\operatorname{KL}(\rho\|\pi)}{n+1}.
\]
Consequently the infimum over all such comparison laws may be taken on
the right-hand side. The theorem requires neither convexity of the
support of $\pi$ nor a Bernstein condition relative to the comparator.
\end{theorem}

\begin{proof}
We first prove an averaged one-observation inequality. Fix a probability
law $q$ on $U\in[-1,1]$, with mean $\mu$ and variance $v$. For a real
label $y$, let $L_y(u)=(u-y)^2$ and
\[
 dp_y(u)=\frac{e^{-L_y(u)/T}}{\int e^{-L_y/T}\,dq}\,dq(u).
\]
For two independent draws $U,U'$ from $q$, the covariance identity gives
\begin{align*}
 \int L_y\,dq-\int L_y\,dp_y
 &=\frac{\mathbb E[(L_y(U)-L_y(U'))
         (e^{-L_y(U')/T}-e^{-L_y(U)/T})]}{2\int e^{-L_y/T}\,dq}.
\end{align*}
Since
\[
 |L_y(u)-L_y(u')|\le2(1+|y|)|u-u'|,
 \qquad \operatorname{osc}_{[-1,1]}L_y\le4(1+|y|),
\]
the mean-value theorem bounds the last display by
\[
 \frac{4(1+|y|)^2}{T}
       e^{4(1+|y|)/T}v.
\]
Also $\int L_y\,dq=L_y(\mu)+v$. Hence
\begin{equation}
 \int L_y\,dp_y\ge L_y(\mu)+
 \left[1-\frac{4(1+|y|)^2}{T}e^{4(1+|y|)/T}\right]v.
 \label{eq:one-step-pointwise}
\end{equation}

Now let $Y=a+\xi$, where $|a|\le1$ and $\xi$ is centered
$\sigma^2$-sub-Gaussian, independently of the fixed law $q$.
The elementary estimates
\[
 \mathbb E e^{t|\xi|}\le2e^{\sigma^2t^2/2},
 \qquad \mathbb E\xi^4\le16\sigma^4
\]
and Cauchy--Schwarz imply
\begin{align*}
 &\mathbb E\left[(1+|Y|)^2e^{4(1+|Y|)/T}\right]\\
 &\quad\le e^{8/T}
       \mathbb E[(8+2\xi^2)e^{4|\xi|/T}]\\
 &\quad\le e^{8/T}
       \left[16e^{8\sigma^2/T^2}
              +8\sqrt2\,\sigma^2e^{16\sigma^2/T^2}\right]
 \le32s
 \qquad(T\ge256s).
\end{align*}
Thus the expected coefficient subtracted in
\eqref{eq:one-step-pointwise} is at most $128s/T\le1/2$.
After averaging over $Y$ we obtain
\begin{equation}
 \mathbb E_Y L_Y(\mu)
 \le\mathbb E_Y\int L_Y\,dp_Y-\tfrac12v
 \le\mathbb E_Y\int L_Y\,dp_Y.
 \label{eq:averaged-one-step}
\end{equation}
The same statement holds conditionally on any input $X=x$ and on an
independent training sample.

Introduce $n+1$ IID observations, retaining the same fixed architecture,
prior and $T$ used by the $n$-observation learner. Let $Q^+$ be their full
posterior and $Q^{-i}$ the posterior after removing observation $i$.
Conditional on the other observations and $X_i$, push $Q^{-i}$ forward
through $w\mapsto f_w(X_i)$ to obtain $q$. It is independent of $Y_i$
given this conditioning. Updating $q$ by the omitted squared-loss factor
gives exactly the corresponding pushforward of $Q^+$.
Applying \eqref{eq:averaged-one-step} and then averaging yields
\[
 \mathbb E (\bar f^{-i}(X_i)-Y_i)^2
 \le\mathbb E\int(f_w(X_i)-Y_i)^2\,dQ^+(w).
\]
Sum over $i=1,\ldots,n+1$. By exchangeability and independence, the
left-hand average is $\mathbb E_{S_n}R(\bar f_{S_n})$.
For each realized enlarged sample the Gibbs variational inequality gives
\begin{align*}
 &\int\sum_{i=1}^{n+1}(f_w(X_i)-Y_i)^2\,dQ^+(w)\\
 &\quad\le\int\sum_{i=1}^{n+1}(f_w(X_i)-Y_i)^2\,d\rho(w)
             +T\operatorname{KL}(\rho\|\pi),
\end{align*}
where the nonnegative term $T\operatorname{KL}(Q^+\|\pi)$ was dropped.
Take expectations and divide by $n+1$ to prove the theorem.
All integrals are finite: the outputs are bounded and the labels have
finite second moment. The posterior normalizers are positive, and the
loss oscillation over parameters is finite for every realized sample.
\end{proof}

\paragraph{Attribution.}
The leave-one-out strategy and a pointwise inverse-tilting inequality
for bounded losses are from M.~M.~H\o gsgaard, P.~Rebeschini and
T.~Wegel, \emph{Aggregation with Exponential Weights is Optimal in
Expectation}, arXiv:2607.02247, July 2026~\cite{hogsgaard}. Their main statement uses a
finite uniform dictionary and bounded losses. Above, sums are replaced
by integrals and the variational comparison is written in KL form;
the pointwise bounded-loss step is replaced by an averaged cooling
inequality for bounded predictions and conditionally sub-Gaussian
labels. The arbitrary-prior and conditional sub-Gaussian extension is proved here.

\section{An optimistic oracle for one final posterior}
\label{sec:statistics}
Throughout this section labels lie in $[-1,1]$. The argument uses only bounded loss and the comparison construction. For every fixed teacher $g\in\calG_{m,B}$ and every $0<r\le1$, Lemma~\ref{lem:comparison} supplies a comparison law with
\begin{equation}
 \E_\nu\whR_S(f_w)\le\whR_S(g)+r,
 \qquad \KL(\nu\Vert\mu_N)\le K_{m,B}r^{-q},
 \qquad q=\alpha/(2-\alpha).
 \label{eq:stat-comparison}
\end{equation}
The law may depend on $g$ and $r$ but not on the observed labels. The argument below also allows sample-dependent comparison laws if these two inequalities hold; the PAC--Bayes event is simultaneous in all posterior laws.

Write
\[
 F_S(\beta)=-\frac1{\beta n}\log Z_{S,\beta}.
\]
For $n\ge2$, put $J=\lceil\log_2(64n)\rceil$ and use the finite grid
$\mathcal B_n=\{2^{-j}:3\le j\le J\}$. Let $M_n=|\mathcal B_n|$,
$\ell_n=\log(2M_n/\delta)$, and define
\begin{equation}
 c_\beta=\frac{1-e^{-4\beta}}{4\beta},\qquad
 \mathcal C_S(\beta)=
 \frac{F_S(\beta)+\ell_n/(\beta n)}{c_\beta}.
 \label{eq:multiplicative-certificate}
\end{equation}
Choose a grid value minimizing $\mathcal C_S$, breaking ties by the smaller index. This is a measurable, finite selection using only the data and the learner's partition function. It selects one likelihood and one posterior, without averaging predictions across temperatures.

\begin{lemma}[A simultaneous bounded-loss certificate]
\label{lem:multiplicative-PAC}
For an architecture fixed before observing the data, with probability at least $1-\delta/2$, simultaneously for $\beta\in\mathcal B_n$,
\[
 R(\bar f_{S,\beta})\le\E_{\Pi_{S,\beta}}R(f_w)
                   \le\mathcal C_S(\beta).
\]
\end{lemma}
\begin{proof}
For every $0\le l\le4$, convexity of $e^{-\beta l}$ gives
\[
 e^{-\beta l}\le1-\frac{1-e^{-4\beta}}4l.
\]
Thus, for each fixed $w$,
$\E e^{-\beta(f_w(X)-Y)^2}\le e^{-\beta c_\beta R(f_w)}$.
Independence implies
\[
 \E_S\exp\{\beta n[c_\beta R(f_w)-\whR_S(f_w)]\}\le1.
\]
Integrate over the prior, use Markov's inequality, and then use the relative-entropy variational inequality. With failure probability at most $e^{-\ell_n}$, simultaneously for every probability law $\rho$,
\[
 c_\beta\E_\rho R(f_w)\le\E_\rho\whR_S(f_w)
       +\frac{\KL(\rho\Vert\mu_N)+\ell_n}{\beta n}.
\]
A union bound over the grid gives the declared event. At its Gibbs posterior the empirical loss and entropy terms sum exactly to $F_S(\beta)$. Jensen's inequality for squared loss proves the ordinary posterior-mean bound. All losses are bounded, so every integral used here is finite.
\end{proof}

\begin{theorem}[Risk-adaptive comparison bound]
\label{thm:optimistic-statistics}
Fix $g\in\calG_{m,B}$ and write $T=K_{m,B}+\ell_n$. With probability at least $1-\delta$, the selected posterior satisfies
\begin{equation}
 R(\bar f_{S,\widehat\beta})\le R(g)+C_q\left[
       \left(\frac{T R(g)}n\right)^{1/(q+2)}
       +\left(\frac Tn\right)^{1/(q+1)}\right].
 \label{eq:optimistic-statistical-oracle}
\end{equation}
The same bound holds for its posterior-average risk. The constant is independent of learner width. Neither the temperature selection nor its certificate uses $g$, $m$, $B$, $K_{m,B}$, or $R(g)$.
\end{theorem}
\begin{proof}
Put $R=R(g)$ and
\[
 u=\left(\frac{TR}n\right)^{1/(q+2)}
       +\left(\frac Tn\right)^{1/(q+1)}.
\]
If $u>1$, the result follows from risk at most four after increasing the constant. Suppose $u\le1$. Define the analytical comparison temperature
$\beta_*=u/[8(R+u)]$. Since $0\le R\le4$ and $T\ge1$,
\[
 \frac1{40n}\le\frac{u}{40}\le\beta_*\le\frac18.
\]
The first inequality uses $u\ge(T/n)^{1/(q+1)}\ge1/n$.
The grid therefore contains a $\beta\in[\beta_*/2,\beta_*]$.
For $\beta\le1/8$, Taylor's inequality gives
$c_\beta\ge1-2\beta$, hence $c_\beta^{-1}\le1+3\beta$.

Use \eqref{eq:stat-comparison} at precision $r=u$ in the free-energy variational formula. Since $u\le1$,
\[
 F_S(\beta)+\frac{\ell_n}{\beta n}
 \le\whR_S(g)+u+\frac{T}{\beta n u^q}.
\]
The definition of $u$ gives
$nu^{q+2}\ge TR$ and $nu^{q+1}\ge T$. Thus
\[
 \frac{T}{\beta n u^q}
 \le\frac{16T(R+u)}{nu^{q+1}}\le32u,
 \qquad \beta R\le u/8.
\]
On the certificate event, the selected certificate is no larger than this grid certificate. Consequently
\[
 \E_{\Pi_{S,\widehat\beta}}R(f_w)
 \le (1+3\beta)\{\whR_S(g)+33u\}.
\]
For the single fixed teacher, bounded-loss Bernstein concentration gives, on a further event of probability at least $1-\delta/2$,
\[
 \whR_S(g)\le R+C\sqrt{R\ell_n/n}+C\ell_n/n.
\]
Both additional terms are $O(u)$. Indeed $u\le1$ implies $T/n\le1$; the power $1/(q+2)<1/2$ makes the first summand in $u$ dominate $\sqrt{RT/n}$, and the power $1/(q+1)<1$ makes the second dominate $T/n$. Substitute this inequality and use $\beta R\le u/8$. The leading population-risk term remains exactly $R$; all remaining contributions are $O_q(u)$. Jensen gives the mean-prediction claim.
\end{proof}

The bounded coefficient class is compact, and its bounded-output risk is continuous in the parameters. Therefore a deterministic population-risk minimizer exists and may be used in Theorem~\ref{thm:optimistic-statistics}. More generally a deterministic sequence approaching the infimum gives the same conclusion. Summable confidence allocation over integer teacher widths and integer coefficient caps yields a simultaneous structural oracle, with the additional logarithms absorbed in $T$. The PAC--Bayes certificate itself remains common to all those comparison classes.

For the Student tail index,
\[
 \frac1{q+2}=\frac{2-\alpha}{4-\alpha},\qquad
 \frac1{q+1}=\frac{2-\alpha}{2}.
\]
Thus a sufficient sample budget for excess at most $\varepsilon$ is
\[
 n\ge\widetilde O_{\alpha}\left(
 K_{m,B}\{R_*\varepsilon^{-(q+2)}
                       +\varepsilon^{-(q+1)}\}\right),
 \qquad R_*=\inf_{g\in\calG_{m,B}}R(g).
\]
For Cauchy weights this gives the oracle remainder
$\widetilde O((K_{m,B}R_*/n)^{1/3}+(K_{m,B}/n)^{1/2})$.
For $\alpha=1/2$ the two powers are $3/7$ and $3/4$.
The architecture is fixed before the observations. These statements have constants uniform in learner width, but do not assert one common PAC--Bayes event for data-dependent selection among infinitely many architectures.

\subsection{A faster expected oracle for averaged prefix predictions}
The ordinary final posterior above can be distinguished from a second, explicitly different predictor. Fix $\eta=1/8$, keep the same architecture and prior, and form the posterior after each prefix of the data with likelihood $\exp(-\eta\sum l_i)$. Let $p_i(x)$ be its ordinary mean after $i$ observations, including $p_0$ from the prior. Define
\[
 \widetilde f_n(x)=\frac1n\sum_{i=0}^{n-1}p_i(x).
\]
This averages prefix predictions. It is not the mean of the posterior after all $n$ observations, and it is not the temperature selection used earlier. The Gaussian observation variance is the constant $1/(2\eta)=4$, independent of sample size and learner width.

\begin{corollary}[Expected oracle for progressive averaging]
\label{cor:progressive-average}
Under the same comparison bound, the averaged predictor satisfies
\[
 \E_S R(\widetilde f_n)
 \le\inf_{g\in\calG_{m,B}}R(g)
       +C_q\left(\frac{K_{m,B}}n\right)^{1/(q+1)}.
\]
For Cauchy weights the power is $1/2$; for $\alpha=1/2$ it is $3/4$. The expectation is over the IID training sample, and constants remain uniform over every architecture fixed in advance with $N\ge m$.
\end{corollary}
\begin{proof}
For $|y|\le1$, the function $a\mapsto e^{-\eta(a-y)^2}$ is concave on $[-1,1]$ when $\eta\le1/8$. Jensen's inequality at step $i$ therefore gives
\[
 (p_{i-1}(X_i)-Y_i)^2\le
 -\frac1\eta\log\E_{\Pi_{i-1}}e^{-\eta(f_w(X_i)-Y_i)^2}.
\]
The right sides telescope through the successive normalizers. The relative-entropy variational formula consequently yields the deterministic inequality
\[
 \sum_{i=1}^n(p_{i-1}(X_i)-Y_i)^2
 \le\inf_\rho\left\{\E_\rho\sum_{i=1}^n(f_w(X_i)-Y_i)^2
                          +\frac{\KL(\rho\Vert\mu_N)}\eta\right\}.
\]
Insert the comparison law at precision $r$ from \eqref{eq:stat-comparison}. The right side is at most
$n\whR_S(g)+nr+K_{m,B}r^{-q}/\eta$.
By independence of the next example from its preceding prefix,
\[
 \E_S\frac1n\sum_{i=1}^n(p_{i-1}(X_i)-Y_i)^2
     =\E_S\frac1n\sum_{i=0}^{n-1}R(p_i).
\]
Convexity of squared risk bounds $R(\widetilde f_n)$ by the average on the right. Taking expectations and optimizing $0<r\le1$ gives the stated rate. If the optimizing scale exceeds one, bounded risk gives the same inequality after increasing $C_q$.
\end{proof}

A direct high-probability version follows by a bounded martingale concentration bound for the prefix losses and scalar concentration for the fixed comparator. It adds $O(\sqrt{\log(1/\delta)/n})$ to the displayed oracle. Thus the faster-than-$n^{-1/2}$ powers here are asserted for expected population risk, not for that elementary high-probability conversion.

\section{A function prior and total online mistakes on finite domains}
\label{sec:uniform-online}
The same independent-coordinate prior gives a stronger kind of guarantee on a finite input domain: a margin-implementable Boolean function receives substantial prior probability, uniformly over learner width. Bayesian sign conditioning then makes only polynomially many mistakes on every input sequence. For this section, the hidden activation need only be bounded by one and $A$-Lipschitz; the second-derivative assumption is unnecessary. Prediction uses the sign of the raw readout.

Let $\mathcal X\subset[-1,1]^d$ be nonempty and finite. Fix a width-$m$, depth-$H$ teacher with coefficients bounded by $B\ge1$, and an accuracy $0<t\le1$. Define
\begin{align}
 M&=(m+1)^H,&G&=1+AmB,& U&=G^H/t,\\
 \Lambda&=\log(4M|\mathcal X|),&
 D_0&=d+1+\Lambda+\sqrt{m\Lambda},\\
 \epsilon&=\min\left\{\frac1{C_\alpha D_0U},
              \left(\frac1{C_\alpha\Lambda U^2}\right)^{1/(2-\alpha)}\right\}.
 \label{eq:uniform-epsilon}
\end{align}
The constant $C_\alpha$ is sufficiently large. Put
\begin{equation}
 Q=C_\alpha M\left[\epsilon^{-\alpha}+(d+1)
 \log\frac{C_\alpha M(d+1)(B+6)^{1+\alpha}}\epsilon\right]+2\log2.
 \label{eq:uniform-resource}
\end{equation}
At fixed depth and tail index, $Q$ is polynomial in $d,m,B,1/t$ and $\log|\mathcal X|$, up to logarithms, and contains no learner-width factor.

\begin{theorem}[Uniform finite-domain prior mass]\label{thm:uniform-mass}
For every finite learner width $N\ge m$,
\begin{equation}
 \mu_N\left\{\max_{x\in\mathcal X}|z_w(x)-z_g(x)|\le t\right\}
 \ge e^{-Q}.
 \label{eq:uniform-mass}
\end{equation}
\end{theorem}
\begin{proof}
Use the comparison measure from Section~\ref{sec:comparison} at precision~\eqref{eq:uniform-epsilon}, and condition first on a fixed index assignment. At a selected row beyond the first, the centered perturbation noise at input $x$ is
\[
 Z(x)=\sum_{\mathrm{designated}}(W_j-a_j)h_j(x)
        +\sum_{\mathrm{background}}W_jh_j(x).
\]
Conditional on the preceding layers, these summands are independent and centered. They have absolute value at most $9\epsilon$: for a small designated coefficient, $|a_j|\le4\epsilon$ and $|W_j|\le5\epsilon$; the other cases are smaller. Their total variance is at most
\[
 V=17m\epsilon^2+C_\alpha\epsilon^{2-\alpha}.
\]
Conditional Bernstein therefore gives, with $u=1/U$,
\[
 \Pp\{|Z(x)|>u\mid\text{preceding layers}\}
 \le2\exp\left\{-\frac{u^2}{2(V+3\epsilon u)}\right\}.
\]
The cutoff~\eqref{eq:uniform-epsilon} ensures $V\le u^2/(4\Lambda)$, $3\epsilon u\le u^2/(4\Lambda)$, and $(d+1)\epsilon\le u$. There are at most $M$ selected later rows. A union bound over them and all input points bounds the probability of any failure by $2M|\mathcal X|e^{-\Lambda}=1/2$. Conditional bounds suffice here; different input evaluations need not be independent.

On the common success event, first-layer preactivation errors are deterministically at most $(d+1)\epsilon\le u$. Let $e_\ell$ be the largest selected hidden-coordinate error over $\mathcal X$. With $a=AmB$,
\[
 e_1\le Au,\qquad e_\ell\le ae_{\ell-1}+Au,\qquad
 \max_{\mathcal X}|z_w-z_g|\le mB e_H+u
 \le u\sum_{j=0}^H a^j\le u(1+a)^H=t.
\]
The success probability is at least $1/2$ for every embedding, hence also for the comparison mixture. Its KL bound~\eqref{eq:embedding-kl} and binary relative entropy give~\eqref{eq:uniform-mass}, exactly as in Corollary~\ref{cor:mass}.
\end{proof}

For Cauchy weights and $\mathcal X=\{\pm1\}^d$, one can simplify~\eqref{eq:uniform-resource} to
\begin{equation}
 Q=\widetilde O_H\left(
 \frac{(d+1)(m+1)^H(1+AmB)^{2H}}{t^2}\right).
 \label{eq:boolean-cauchy-resource}
\end{equation}
Indeed $\Lambda=O(d+H\log(m+1))$ and $U\ge m^H\ge\sqrt m$, so the terms $D_0U$ are absorbed by a constant times $(d+H\log(m+1))U^2$. No exponential in input dimension is suppressed in~\eqref{eq:boolean-cauchy-resource}.

\begin{corollary}[Prior mass of a Boolean function]\label{cor:sign-prior}
Suppose $f:\mathcal X\to\{\pm1\}$ has a teacher satisfying
$f(x)z_g(x)\ge2\kappa$ at every input, where $\kappa>0$.
Use $t=\min\{1,\kappa/2\}$ in~\eqref{eq:uniform-resource}. Then
\[
 \mu_N\{\operatorname{sign}z_w=f,\ \min_{\mathcal X}|z_w|>\kappa\}
 \ge e^{-Q}.
\]
The same lower bound holds for the induced sign-function prior if the weight prior is initially conditioned on $\min_{\mathcal X}|z_w|>\kappa$.
\end{corollary}
\begin{proof}
On the uniform-approximation event, $f(x)z_w(x)\ge2\kappa-t>\kappa$. The conditioning event has positive prior mass, and dividing by its probability can only increase the lower bound.
\end{proof}

\begin{corollary}[Every-sequence online guarantees]\label{cor:online-mistakes}
Use either sign-function prior from Corollary~\ref{cor:sign-prior}, and condition exactly on each revealed label $f(X_t)$. Inputs may be chosen adaptively from previous observations and predictions; in the randomized-prediction case, the next input is chosen before drawing the next prediction's fresh randomness. Let $p_t$ be the posterior probability of the correct sign before its reveal. On every realized input sequence,
\begin{equation}
 \sum_{t\ge1}-\log p_t\le Q.
 \label{eq:online-logloss}
\end{equation}
Consequently:
\begin{itemize}
\item Fresh posterior-sampled predictions have expected total mistake count at most $Q$.
\item Posterior majority prediction makes at most $Q/\log2$ mistakes, allowing ties to count as mistakes.
\item For every $a\in(0,1)$, there are at most $Q/(-\log a)$ steps at which $p_t\le a$. In particular, failure of strict $>2/3$ correct-label confidence occurs at most $Q/\log(3/2)$ times.
\end{itemize}
These statements require neither IID inputs nor a stationary train--test distribution.
\end{corollary}
\begin{proof}
Let $\pi$ be the induced function prior. For every finite realized sequence, the posterior normalizers telescope:
\[
 \prod_{t=1}^T p_t
 =\pi\{h:h(X_t)=f(X_t)\text{ for every }t\le T\}
 \ge\pi(f)\ge e^{-Q}.
\]
Taking logarithms and then increasing $T$ proves~\eqref{eq:online-logloss}. This is a pathwise identity even when the chosen inputs depend on previous predictions. Since $1-p\le-\log p$, the sum of conditional sampled-prediction error probabilities is at most $Q$; the tower property and monotone convergence give the expected total count. Each majority mistake has $p_t\le1/2$, and each confidence failure has $p_t\le a$, proving the other assertions.
\end{proof}

This online result uses noiseless sign observations. It is separate from the Gaussian-likelihood agnostic theorem: the prior is the same, but the observation model and prediction criterion differ. It supplies a literal function-prior lower bound and the corresponding finite total mistake guarantee, rather than only an IID sample-complexity conclusion.

\section{Other heavy-tailed densities and an activation obstruction}
\label{sec:extensions}
\subsection{Bounded Lipschitz activations without smoothness}
\begin{corollary}[The bounded nonsmooth case]\label{cor:bounded-lipschitz}
Suppose the hidden activation is bounded by one and globally Lipschitz, without assuming two derivatives. Let the output link be any fixed bounded Lipschitz map into $[-1,1]$. At the same fixed temperature as Theorem~\ref{thm:fast-oracle}, with probability at least $1-\delta$,
\[
 R(\bar f_{S,\beta})\le R_0+(1+\eta)A_*
 +\frac C\eta\left(
 \frac{(1+\sigma^2)\{K^{\mathrm{Lip}}_{m,B}+\log(2/\delta)\}}n
 \right)^{(2-\alpha)/2}.
\]
The resource $K^{\mathrm{Lip}}_{m,B}$ has the same polynomial orders in $d,m,B$ as~\eqref{eq:complexity}, with constants adjusted for the activation and output link. Smoothness improves the dependence on $\eta$ in the main theorem, but is unnecessary for its sample-size power.
\end{corollary}
\begin{proof}
The variance induction in Lemma~\ref{lem:perturbation} uses only Lipschitzness and bounded hidden features. Thus, after a Lipschitz output link, a comparison law can satisfy
\[
 \sup_x\E_\nu|f_w(x)-g(x)|^2\le r,\qquad
 \KL(\nu\|\mu_N)\le K^{\mathrm{Lip}}_{m,B}r^{-q},
 \qquad q=\frac\alpha{2-\alpha},
\]
by choosing $\epsilon$ proportional to $(r/S_{m,B})^{1/(2-\alpha)}$.
Use the bounded regression reference $f_0$ and the squared triangle inequality:
\[
 a_\nu:=\E_\nu A(f_w)
 \le(1+\eta/8)A(g)+(1+8/\eta)r.
\]
The relative-loss PAC--Bayes and comparison-concentration lemmas then give the same proof as Theorem~\ref{thm:fast-oracle}, except that its leading approximation factor is bounded by
$(1+\eta/8)^2/(1-\eta/8)\le1+\eta$ and the remainder is at most
\[
 \frac C\eta\left[r+\frac{(1+\sigma^2)
              \{K^{\mathrm{Lip}}_{m,B}r^{-q}+\log(2/\delta)\}}n\right].
\]
Optimize at $r=[(1+\sigma^2)(K^{\mathrm{Lip}}_{m,B}+\log(2/\delta))/n]^{1/(q+1)}$, using bounded risk when this exceeds one.
\end{proof}

\begin{corollary}[A sharp expectation bound without smoothness]
In the bounded-Lipschitz setting above, put $q=\alpha/(2-\alpha)$,
$u=(1+\sigma^2)K^{\rm Lip}_{m,B}/(n+1)$ and $A_g=R(g)-R_0$.
At the fixed temperature $\beta_*=1/[256(1+\sigma^2)]$,
\[
 \E R(\bar f_{S_n,\beta_*})
 \le R(g)+C\{A_g^{q/(2q+1)}u^{1/(2q+1)}
                         +u^{1/(q+1)}\}.
\]
The infimum over deterministic comparators can be taken directly.
\end{corollary}
\begin{proof}
The mean-square comparison gives risk at most
$R(g)+2\sqrt{A_g r}+r$ and KL at most $K^{\rm Lip}_{m,B}r^{-q}$.
Apply Theorem~\ref{thm:generic-expectation}. For $u\le1$, choose $r$ as the smaller of $u^{1/(q+1)}$ and $(u/\sqrt{A_g})^{1/(q+1/2)}$, omitting the latter when $A_g=0$. These choices bound the approximation and KL terms by the two displayed powers. For $u>1$, bounded excess risk suffices.
\end{proof}

\subsection{Comparable densities and symmetric stable laws}
The Student formula is a convenient explicit choice, rather than an essential part of the argument.

\begin{corollary}[A class of independent-coordinate priors]\label{cor:comparable-density}
All preceding conclusions remain valid if the base Student density is replaced by any symmetric density $p$ satisfying
\begin{equation}
 c_-(1+t^2)^{-(1+\alpha)/2}
 \le p(t)\le c_+(1+t^2)^{-(1+\alpha)/2}
 \quad(t\in\R),
 \label{eq:comparable-density}
\end{equation}
for fixed $0<c_-\le c_+<\infty$. The constants may additionally depend on $c_-,c_+$. Coordinates remain independent and a row of fan-in $F$ still uses scale $F^{-1/\alpha}$.
\end{corollary}
\begin{proof}
The proof of Lemma~\ref{lem:truncation} uses only symmetry, a positive lower density near zero, the upper tail bound $C u^{-\alpha}$, the truncated second-moment bound $C u^{2-\alpha}$, and the lower density bound~\eqref{eq:density-lower}. Each follows from~\eqref{eq:comparable-density}, with modified constants. All later steps depend only on those estimates. In particular, one does not compare the full product density to the Student product density, which would incur an unnecessary factor per unused coordinate.
\end{proof}

A familiar example is the symmetric $\alpha$-stable law with characteristic function $\E e^{itT}=e^{-|t|^\alpha}$. Its density is positive and continuous, and its tail satisfies
\[
 p(t)\sim\frac{\Gamma(1+\alpha)\sin(\pi\alpha/2)}{\pi|t|^{1+\alpha}}
 \quad(|t|\longrightarrow\infty);
\]
these standard stable-density facts imply~\eqref{eq:comparable-density} \cite{zolotarev}.
For this family the fan-in scale is particularly natural: a sum of $F$ independent copies scaled by $F^{-1/\alpha}$ has exactly the original stable law. More generally, conditional on bounded hidden values $h_j$, the next preactivation is stable with scale
\[
 \left(\frac1F\sum_{j=1}^F|h_j|^\alpha\right)^{1/\alpha}\le1.
\]
Thus bounded hidden activations also keep the native layer scales controlled.

\subsection{Nonvanishing initialization scale and small-network prior mass}
\label{sec:cauchy-native-scale}
A nonzero width limit of the initialization output scale does not by itself imply Gaussian-process or kernel behavior. The bounded-activation Cauchy model gives an explicit counterexample to that inference.

\begin{proposition}[The fixed-depth native Cauchy scale]
\label{prop:cauchy-native-scale}
Fix an input $x\in[-1,1]^d$ and hidden depth $H\ge1$. Use tanh hidden activations, a raw linear readout followed by either tanh or clipping to $[-1,1]$, and the independent Cauchy prior with scale $1/(d+1)$ in the first layer and $1/N$ in all subsequent layers, including the readout. Define
\[
 a_0=\frac{1+\sum_{j=1}^d|x_j|}{d+1}>0,
 \qquad \Psi(a)=\E|\tanh(aZ)|,\quad Z\sim\operatorname{Cauchy}(1),
 \qquad a_\ell=\Psi(a_{\ell-1}).
\]
At fixed $H$, as $N\to\infty$,
\[
 a_{H,N}:=\frac1N\sum_{j=1}^N|h_j^H(x)|
 \longrightarrow a_H\in(0,1)
 \quad\text{in probability},
 \qquad z_w(x)\ \Longrightarrow\ \operatorname{Cauchy}(a_H).
\]
Conditional on the hidden network, the raw readout has exactly the Cauchy law with scale $a_{H,N}$. In particular its limiting absolute median is $a_H>0$. The raw readout has infinite variance. The actual bounded output, however, has finite variance satisfying
\[
 \operatorname{Var}(f_w(x))\longrightarrow
 \E[\psi(a_H Z)^2]\in(0,1).
\]
Thus its initialization standard deviation is bounded above and below by positive constants as width increases at fixed depth.
\end{proposition}
\begin{proof}
The characteristic function of a centered Cauchy variable of scale $b$ is $e^{-b|t|}$. Independence therefore shows that $\sum_j c_jZ_j$, for independent scale-$b_j$ Cauchy variables, is Cauchy of scale $\sum_j|c_j|b_j$.

First-layer preactivations at the fixed input are consequently independent Cauchy variables of scale $a_0$. Their bounded transformed absolute values have empirical average converging to $a_1$. At every later hidden layer, conditional on the preceding hidden vector, the preactivations are independent Cauchy variables of scale $a_{\ell-1,N}$. Their transformed absolute values lie in $[0,1]$, so conditional Hoeffding gives, for every $u>0$,
\[
 \Pr\{|a_{\ell,N}-\Psi(a_{\ell-1,N})|>u
                 \mid h^{\ell-1}(x)\}\le2e^{-2Nu^2}.
\]
The function $\Psi$ is continuous on $[0,1]$ by dominated convergence, and is strictly positive and less than one at every positive argument. Induction in the fixed number of layers proves $a_{H,N}\to a_H\in(0,1)$ in probability.

The same Cauchy addition identity gives the exact conditional readout law. Its unconditional characteristic function is $\E e^{-|t|a_{H,N}}$, which converges to $e^{-|t|a_H}$ by bounded convergence in probability. This proves weak convergence. The limiting absolute-value distribution is continuous and strictly increasing, with median $a_H$, so the absolute medians converge as well. At every finite width the first-layer continuous variables give positive hidden absolute averages almost surely; conditioning on any positive readout scale gives an infinite raw second moment. The bounded output $\psi(z)$ is symmetric and therefore has mean zero. Its variance converges by weak convergence and the bounded continuous function $z\mapsto\psi(z)^2$, proving the final assertion.
\end{proof}

The earlier prior-mass and online-learning conclusions apply to this same tanh Cauchy prior. Thus a nonvanishing fixed-depth initialization scale coexists with width-independent small-network neighborhoods and polynomial total online mistakes. The Gaussian central-limit intuition fails because the coordinates have infinite variance and their sums remain stable Cauchy variables. The proposition does not assert a scale lower bound uniform in growing depth.


\subsection{Why unbounded ReLU is different for this prior}
The bounded-hidden assumption has a concrete role. With the same Cauchy fan-in scaling, an unbounded ReLU hidden layer can amplify the raw output scale as width increases. A final tanh or clip alone does not repair the resulting fixed-temperature posterior.

\begin{proposition}[A one-layer Cauchy--ReLU obstruction]\label{prop:relu-obstruction}
Replace the hidden activation by ReLU, keep one hidden layer, a final tanh or clip to $[-1,1]$, and the independent Cauchy prior with scales $1/(d+1)$ and $1/N$. Fix a temperature $\beta>0$. There is a noiseless distribution with a width-one, coefficient-bound-one teacher for which no width-uniform vanishing-risk oracle for the ordinary final posterior mean can hold at this temperature.
\end{proposition}
\begin{proof}
Take $d=1$, the single input $x=0$, and the constant label $c=1/2$. A width-one teacher with readout weight one and first-layer bias $c$ for clipping, or $\operatorname{arctanh}(c)<1$ for a tanh output, produces $c$ exactly. Under the prior, the hidden values are IID
\[
 h_j=\max\{0,B_j\},\qquad B_j\sim\operatorname{Cauchy}(1/2).
\]
They are nonnegative and have infinite mean. Hence $S_N=N^{-1}\sum_{j=1}^Nh_j\to\infty$ almost surely. One way to see this is to apply the strong law to $\min\{h_j,K\}$ for every integer $K$, and then let $K$ increase.

Conditional on the hidden values, stability of the independent readout weights gives
\[
 z_N\mid(h_j)\sim\operatorname{Cauchy}(S_N).
\]
Therefore the distribution of $\psi(z_N)$ converges to equal point masses at $-1$ and $1$. For any fixed $n$, the bounded continuous likelihood weight $e^{-\beta n(f-c)^2}$ can be passed through this limit. The ordinary posterior mean consequently converges, as $N\to\infty$, to
\[
 \frac{e^{-\beta n(1-c)^2}-e^{-\beta n(1+c)^2}}
      {e^{-\beta n(1-c)^2}+e^{-\beta n(1+c)^2}}
 =\tanh(2\beta n c).
\]
Its squared test error tends to $(\tanh(2\beta n c)-c)^2$. Letting $n$ increase after taking this width limit gives $(1-c)^2>0$, whereas the comparator risk is zero. For arbitrarily large sample sizes one can therefore choose a finite learner width with error bounded below by a positive constant. This contradicts any bound tending to zero with $n$ uniformly in width for this fixed-temperature learner.
\end{proof}

This proposition concerns the Cauchy fan-in prior and a fixed likelihood temperature. It does not exclude ReLU guarantees under other priors or different temperature choices. It explains why the bounded-hidden theorem is not automatically an all-activation theorem.


\section{Scope and consequences}\label{sec:biases}
\paragraph{Approximation and noise.}
The main theorem allows arbitrary conditional sub-Gaussian noise around a bounded regression function, including unbounded Gaussian labels. The target need not be a network, and the comparator need not interpolate the observations. Its approximation error is retained up to a freely chosen factor $1+\eta$, while irreducible noise has coefficient exactly one. The expectation theorem has coefficient exactly one on total comparator risk at a fixed temperature. The complementary bounded-label theorem~\eqref{eq:main-oracle} has coefficient one with high probability after calibration.

\paragraph{What is fixed.}
The prior is a single dense product distribution, with independent continuous coordinates. The learner width and depth are fixed before the sample. The comparison proof averages over embeddings of a small computation tree, but that averaging changes neither the prior nor the architecture. The sharp expectation theorem and near-sharp high-probability theorem use fixed temperatures. The separate bounded-label, coefficient-one high-probability theorem selects a temperature from the specified grid; this is its only data-dependent model choice.

\paragraph{Why the output is bounded.}
The raw readout of a Cauchy-weight network need not have a first or second moment. A fixed bounded output makes the posterior mean well defined and controls the relative-loss exponential moments used in the main proof. It also permits the bounded-loss argument for the complementary theorem. The perturbation proof nevertheless controls the raw comparison network in mean square. A semiconcavity property of squared loss after the output link controls the perturbation. For clipping, it also covers teachers whose raw outputs lie exactly on a clipping threshold. For a smooth link such as tanh, the same estimate follows from its bounded derivative and Lipschitz derivative.

\paragraph{Biases beyond the first layer.}
For a nonzero activation, choose a fixed $b_0$ with $c_0=\phi(b_0)\ne0$. A constant hidden wire can be generated from the first-layer constant input and propagated by a coefficient $b_0/c_0$. An affine offset $b$ into a later neuron is then represented by an incoming coefficient $b/c_0$ from that wire. This increases width by one and changes the coefficient cap by a fixed activation-dependent factor. A readout bias is handled in the same way. The theorem therefore includes bounded biased comparators after this explicit compilation, when the learner has room for the added constant wire. If $\phi\equiv0$, this compilation is unavailable; the bias-free statement remains valid.

\paragraph{Depth and activation restrictions.}
Unrolling width $m$ over depth $H$ uses $(m+1)^H$ nodes, and the stability factor contains $(1+AmB)^{2H}$. These are substantial depth costs. The sharper dependence on the approximation factor uses two derivatives. The bounded-Lipschitz extension also gives the same sample-size power for every bounded Lipschitz hidden activation, with a weaker dependence on that factor. Neither statement establishes the same oracle for unbounded ReLU activations under this same-index heavy-tailed prior. The bounded-activation assumption is useful for an additional reason: even unselected background neurons remain uniformly bounded despite heavy-tailed incoming weights.

\paragraph{Verification.}
The probabilistic construction, perturbation calculation and oracle inequality have received independent AI-agent mathematical reviews. This paper has not been formalized in Lean. The accompanying verification directory records the reviewed sources and build.
\begin{thebibliography}{9}
\bibitem{zolotarev} V. M. Zolotarev. \emph{One-dimensional Stable Distributions}. Translations of Mathematical Monographs, vol. 65. American Mathematical Society, 1986.
\bibitem{hogsgaard} M. M. H\o gsgaard, P. Rebeschini, and T. Wegel. Aggregation with Exponential Weights is Optimal in Expectation. arXiv:2607.02247, 2026.
\bibitem{castillo-egels} I. Castillo and P. Egels. Posterior and Variational Inference for Deep Neural Networks with Heavy-Tailed Weights. \emph{Journal of Machine Learning Research} 26(122), 1--58, 2025.
\end{thebibliography}
\end{document}
